\section{Introduction}
Data errors can affect only a few attributes of a record. A suspicious product measurement
does not imply that its name and producer are wrong, and a rare bibliographic entry need
not contain an error. This makes cell localization a different objective from ranking
entire rows as anomalies. Robust cellwise methods already exploit relationships among
variables~\cite{rousseeuw2018ddc}; a foundation model does not make that problem new.
The question addressed here is whether the context used by a conditional predictor can
be improved without access to clean examples or error labels.

A pretrained tabular model offers a practical route to conditional scoring. Treat each
column in turn as a prediction target, use the remaining columns as features, and score
the observed target value under the resulting prediction. TabPFN supports classification
and regression through in-context prediction~\cite{hollmann2025tabpfn}. In a dirty table,
however, the reference examples are themselves imperfect. Repeated corrupt values may
become predictable, while infrequent correct values can receive large suspicion scores.
Removing the highest-scoring context cells' rows could therefore either improve the
reference distribution or erase legitimate variation. Neither outcome follows from
the model's supervised prediction capability.

We study this intervention under a controlled cell-detection protocol. For each target
column and outer query fold, a nested first pass scores only the outer training pool.
The purification arm removes context rows with the largest target-cell scores. A random
arm removes exactly the same number of rows from the same base context. Both retain
the same queries and evaluation denominator, which separates targeted selection from
simply reducing context size; the extra first-pass computation is reported separately.

The empirical scope is deliberately specific: four capped benchmark tables, their
real errors, two controlled corruption conditions and five seeds. Real-error
purification is primary; injections and an observed-key leave-entity-out analysis
are probes. Effects are reported per source table; repeated conditions and seeds
are not treated as additional independent datasets.

The study contributes an outer-fold-isolated implementation of a simple context
intervention, a matched-count and explicitly costed comparison against random removal,
and an artifact-backed assessment spanning accuracy, rare-valid alarms, support,
failures and boundary sensitivity. It does not introduce a new pretrained model,
an influence estimator, or a conformal guarantee. The breadth of any empirical
conclusion is limited to the measured implementations and named tables.
